We estimate the direct effect against continuation at 1.56 percentage points (group-bootstrap 95\% interval [$-$0.12, 3.41]; group sign-flip $p=0.1048$; Holm-adjusted $p=0.2096$). Against the matched arm-outcome comparator we estimate $-$0.39 percentage points (group-bootstrap 95\% interval [$-$2.26, 1.57]; raw group sign-flip $p=0.7927$; Holm-adjusted $p=0.7927$). The two controllers reach 54.43\% and 54.82\%. The denominator contains all 384 planned tasks: 355 reached the frozen checkpoint and 29 ended earlier, and Table~\ref{tab:confirmation-policies} reports the full census. Neither primary controller comparison is significant before or after Holm correction. The development-selected safe policy changes success by 1.95 percentage points relative to continuation. These tests answer the frozen comparisons; a favorable unadjusted subset is not substituted for them.

Direct selects continuation on 245/355 eligible prefixes, warning on 4, replanning on 106, and rollback on 0. Matched selects rollback on 71 eligible prefixes. Direct records 26 draw-level recoveries and 14 disruptions. Matched intervenes on 78.0\% of prefixes, with 66 recoveries and 51 disruptions. Best fixed replanning reaches 55.34\% versus 54.43\% for Direct. Its 74 recoveries and 55 disruptions exceed Direct's counts by 48 and 41, respectively, leaving 7 additional successful draw-level contrasts (0.91 percentage points). These descriptive counts do not identify why the fixed rule scores higher. Direct uses 27.4 calls, 42560 input tokens, and 63.4 client-seconds per planned task, versus 28.1, 43557, and 64.9 for continuation. These summaries are descriptive. Section~\ref{sec:analysis} reads the same branches for what this decision point and this menu could have delivered, which is what the second confirmation then changes.

